\section*{Linear Least-Squares Compositional Fitting}
\color{red}
The phase assignment method described above excludes \qty{23.9}{\percent} (\num{272} of \num{1140}) of the measured indents from the analysis because they fall into grid cells with no dominant phase (>\qty{50}{\percent}).
This exclusion limits the statistical reliability, particularly for rare phases such as \acs{C3A}/\acs{C4AF}.
To address this limitation, a linear least-squares fitting approach was employed to estimate the intrinsic mechanical properties of each phase from the complete dataset.

The mechanical properties of a composite material can be approximated using a linear rule-of-mixtures model:
\begin{equation}
	\label{eq:compositional-fit}
	M_{\text{cell}} = \sum_{i=1}^{n} w_{i,\text{cell}} \cdot M_i
\end{equation}
where $M_{\text{cell}}$ is the measured mechanical property (hardness $H$ or modulus $E_r$) in a grid cell, $w_{i,\text{cell}}$ is the volume fraction of phase $i$ in that cell (determined from the \ac{EDX} grid analysis), and $M_i$ is the intrinsic mechanical property of pure phase $i$ (to be determined).
The cell-averaged properties represent the spatial averaging over the indent interaction volume, which typically extends \qty{\sim 500}{\nano\meter} into the material.

For all \num{1140} grid cells, a system of linear equations for hardness and reduced modulus datasets were constructed and solved using least-squares minimization.
The phase fractions $w_{i,\text{cell}}$ serve as the design matrix (coefficients), and the measured indentation properties $M_{\text{cell}}$ form the observation vector.
The fit parameters $M_i$ were determined by minimizing the residual sum of squares.

Two constraints were applied:
\begin{enumerate}
	\item Pores were fixed to $M_i = 0$ (hardness and modulus; physically reasonable assumption).
	\item Negative property values were eliminated during post-processing (non-physical).
\end{enumerate}

The quality of the fit was assessed using the coefficient of determination ($R^2$).
Uncertainty estimates for the fitted parameters were derived from the covariance matrix of the least-squares solution, calculated as $\text{Cov} = \sigma_{\text{residual}}^2 \cdot (X^T X)^{-1}$, where $X$ is the design matrix and $\sigma_{\text{residual}}^2$ is the residual variance.

For hardness, the fit yields $R^2 = 0.7992$ for the coarse solution and $R^2 = 0.7209$ for the weighted solution, indicating that the linear combination of phase fractions explains the majority of the measured cell-averaged hardness variance.
The fitted phase hardnesses are presented in \zcref{tab:ni-results-fitted}.
The clinker phases (alite/belite: coarse \qty{8.77}{\giga\pascal}, weighted \qty{7.92}{\giga\pascal}; slag: coarse \qty{7.10}{\giga\pascal}, weighted \qty{6.34}{\giga\pascal}) display substantially higher hardness than hydrated phases, consistent with the original observations.
Notably, the fitted value for \acs{C3A}/\acs{C4AF} is coarse \qty{15.10}{\giga\pascal} and weighted \qty{10.46}{\giga\pascal}, derived from all cells containing this phase (even in small fractions), rather than relying on the single direct measurement in the threshold-based approach.
The hydrated phases show fitted hardnesses ranging from coarse \qty{0.79}{\giga\pascal} (C-S-H) to coarse \qty{2.36}{\giga\pascal} (CH), with the weighted fit generally shifting the softer hydrate values slightly upward.
Quartz, with coarse \qty{11.73}{\giga\pascal} and weighted \qty{11.11}{\giga\pascal}, reflects its high hardness.

For the reduced modulus, the compositional fitting approach yields $R^2 = 0.7435$ for the coarse solution and $R^2 = 0.6404$ for the weighted solution, indicating that the linear composition model captures a substantial fraction of the measured modulus variance.
The fitted moduli for clinker phases are alite/belite (coarse \qty{105.51}{\giga\pascal}, weighted \qty{96.09}{\giga\pascal}) and slag (coarse \qty{82.47}{\giga\pascal}, weighted \qty{75.22}{\giga\pascal}).
Hydrated phases show fitted moduli ranging from coarse \qty{20.88}{\giga\pascal} (AFm/AFt) to coarse \qty{50.48}{\giga\pascal} (CH).
C-A-S-H is fitted as coarse \qty{23.88}{\giga\pascal} and weighted \qty{25.80}{\giga\pascal}, C-S-H as coarse \qty{24.01}{\giga\pascal} and weighted \qty{25.63}{\giga\pascal}, and Mg-C-A-S-H as coarse \qty{31.52}{\giga\pascal} in both solutions.
The matrix phase (CSHo equivalent) is fitted as coarse \qty{22.85}{\giga\pascal} and weighted \qty{23.63}{\giga\pascal}.
The \acs{C3A}/\acs{C4AF} phase is estimated as coarse \qty{179.86}{\giga\pascal} and weighted \qty{124.38}{\giga\pascal}, which should be interpreted with caution given the expected sparsity of this phase in the microstructure.
Quartz exhibits the highest fitted modulus at coarse \qty{88.93}{\giga\pascal} and weighted \qty{84.87}{\giga\pascal}, consistent with aggregate behaviour.

\subsubsection*{Advantages and Limitations of the Compositional Fitting Approach}

All \num{1140} measurements are incorporated, eliminating the data loss from the threshold-based method. 
Furthermore, phases like \acs{C3A}/\acs{C4AF}, represented by limited direct measurements, benefit from information distributed across all grid cells, resulting in more stable estimates.
However, the linear rule-of-mixtures assumes additive contributions without interactions.
Real composite behaviour may deviate due to phase-boundary effects, residual stresses, or mechanical constraint effects neighbouring phases.
Additionally, the model assumes each grid cell is internally homogeneous, whereas sub-cellular heterogeneity (micro-cracks, compositional gradients) could introduce errors.
% \textbf{Advantages:}
% \begin{itemize}
% 	\item \textit{Complete data utilisation:} 
% 	\item \textit{Physically constrained solutions:} Fixing pores to zero eliminates unphysical results and reduces the number of free parameters from 11 to 10.
% 	\item \textit{Reduced standard uncertainties:} The incorporation of more datapoints and implicit averaging over phase-composition correlations generally yields lower parameter uncertainties.
% 	\item \textit{Systematic comparison:} Properties derived from the same optimisation criterion are directly comparable.
% \end{itemize}

% \textbf{Limitations and Caveats:}
% \begin{itemize}
% 	\item \textit{Model assumptions:} 
% 	\item \textit{Homogeneity assumption:} The model assumes each grid cell is internally homogeneous, whereas sub-cellular heterogeneity (micro-cracks, compositional gradients) could introduce error.
% 	\item \textit{Interaction volume:} The indentation interaction volume (\qty{\sim 500}{\nano\meter}) is significantly larger than individual grid cells (\qty{\sim 1}{\micro\meter}), potentially averaging across phase boundaries and introducing correlation between measured properties and measured composition.
% 	\item \textit{Spatial correlation:} Fitted parameters may reflect spatial correlations in the microstructure (e.g. certain phases preferentially aggregate) rather than true intrinsic phase properties.
% 	\item \textit{C3A/C4AF and Quartz uncertainties:} These phases are either rare or belong to the aggregate; fitted values may not represent intrinsic single-phase properties.
% \end{itemize}

\subsection*{Comparison: Threshold-Based vs. Compositional Fitting}

The two approaches yield broadly consistent results for abundant phases but differ in approach and reliability for sparse phases.
For C-S-H (the dominant hydrate), the threshold-based method yields \qty{1.06(0.53)}{\giga\pascal} (median \qty{1.00}{\giga\pascal}), while the compositional fit estimates \qty{0.82}{\giga\pascal}.
This difference likely reflects the fitting method's averaging of many weakly-contributing cells against the threshold method's focus on cells dominated by a single phase.
For quartz, both methods agree closely: \qty{11.01(2.29)}{\giga\pascal} vs. \qty{11.75}{\giga\pascal}.

The fitted approach provides a complementary perspective: rather than asking \textit{``what is the hardness of indents in cells dominated by phase X?''}, it answers \textit{``what intrinsic phase property best explains the observed cell-averaged hardness across all observed phase combinations?''}
This distinction is important for phases present in small fractions and for assessing composite-scale mechanical behaviour.


\begin{table}[htb!]
	\centering
	\caption{
		Fitted phase-dependent mechanical properties from the compositional least-squares model. Values are presented for the coarse and weighted fitting variants.
	}
	\footnotesize
	\begin{tabularx}{\linewidth}{
		X
		S[table-format=2.2, round-mode=places, round-precision=2]
		S[table-format=2.2, round-mode=places, round-precision=2]
		S[table-format=3.2, round-mode=places, round-precision=2]
		S[table-format=3.2, round-mode=places, round-precision=2]
	}
		\toprule
		Phase & \multicolumn{2}{c}{Hardness $H$ (\unit{\giga\pascal})} & \multicolumn{2}{c}{Modulus $E_r$ (\unit{\giga\pascal})} \\
		& {Coarse} & {Weighted} & {Coarse} & {Weighted} \\
		\midrule
		Quartz & 11.73 & 11.11 & 88.93 & 84.87 \\
		\acs{C3S}/\acs{C2S} & 8.77 & 7.92 & 105.51 & 96.09 \\
		Slag & 7.10 & 6.34 & 82.47 & 75.22 \\
		\acs{CH} & 2.36 & 2.10 & 50.48 & 42.60 \\
		\acs{CSHi} & 0.79 & 1.06 & 24.01 & 25.63 \\
		\acs{C3A}/\acs{C4AF} & 15.10 & 10.46 & 179.86 & 124.38 \\
		Mg-\ac{CASH} & 1.32 & 1.45 & 31.52 & 31.52 \\
		\ac{CASH} & 0.95 & 1.24 & 23.88 & 25.80 \\
		\acs{AFm}/\ac{AFt} & 1.09 & 1.19 & 20.88 & 25.01 \\
		\acs{CSHo} & 1.05 & 1.07 & 22.85 & 23.63 \\
		Pores & 0.40 & 0.40 & 5.00 & 5.00 \\
		\midrule
		R$^2$ & 0.7992 & 0.7209 & 0.7435 & 0.6404 \\
		\bottomrule
	\end{tabularx}
	\label{tab:ni-results-fitted}
\end{table}

% ---

% \subsection*{Graphical Representations}

% The following figures are recommended to visualise the fitting quality and the fitted parameters:

% \subsubsection*{(a) Measured vs. Predicted Cell Properties}

% A \textbf{scatter plot} of measured cell-averaged hardness (or modulus) vs. predicted values from the compositional fit, colour-coded by dominant phase or overall phase diversity index.
% The diagonal line $y=x$ indicates perfect prediction. Points below the diagonal are underestimated by the model; points above are overestimated.
% This reveals systematic deviations (e.g. interaction-volume boundary effects) and outliers.

% \subsubsection*{(b) Comparison Bar Chart}

% \textbf{Grouped bar charts} showing threshold-based vs. fitted values for each phase, enabling direct visual comparison.
% Error bars on the threshold-based values represent standard deviations; fitted values can include confidence intervals from the covariance matrix.

% \subsubsection*{(c) Residual Analysis}

% A \textbf{histogram or Q-Q plot} of the fit residuals ($M_{\text{measured}} - M_{\text{predicted}}$) to assess normality and detect systematic patterns.
% A heatmap showing residuals as a function of spatial location and composition could reveal whether errors correlate with specific microstructural features.

% \subsubsection*{(d) Phase-Property Space}

% A \textbf{scatter plot} in phase-composition space (e.g. fraction of hydrates vs. measured hardness), with points coloured by the fitted property value.
% Contour lines overlay the fitted linear model, illustrating the compositional relationship.

% \subsubsection*{(e) Sensitivity/Correlation Heatmap}

% A \textbf{heatmap} of the correlation matrix derived from the covariance of fitted parameters, highlighting which phase properties are mechanically coupled via microstructural correlations.
\color{black}